% !PRINTED 273
\weberpage{273}
If we conceive the group $L_p$ as the group of the linear substitutions
\[
\eta = \frac{a\xi+b}{c\xi+d}
\]
and let $\xi, \eta$ in it run through the numbers $\infty, 0, 1, \ldots, p-1$ according to \S.\,65, then $G$ gives us the group of the integral linear substitutions
\[
\eta = a^2\xi + ab.
\]
These are the substitutions that leave the figure $\infty$ unchanged, and hence yield a permutation group of only $p$ figures. These are the same special metacyclic groups that we considered in the seventeenth chapter of the first volume.

The various groups transformed from $G$ leave some one of the remaining $p+1$ indices unchanged, and the total number of these groups is $p+1$.

\begin{center}
\S.\,71.\\
\gesperrt{Divisors of the group $L_p$ whose degree is not divisible}\\
\gesperrt{by $p$.}
\end{center}

To find the remaining divisors of $L_p$ whose degree is not divisible by $p$, we can go, step by step, the same way that led us in the seventh and eighth chapters to the determination of the polyhedral groups. We may content ourselves here with bringing out the principal points of the derivation, and refer for the proofs to the place named, where entirely the same conclusions were to be drawn.

For this it is required to conceive the group $L_p$, as set forth in \S.\,65, as a group of linear fractional substitutions
\weq{1}{\Theta(\xi) = \frac{a\xi+b}{c\xi+d}}
and to assign to the variable $\xi$ in it all $p^2+1$ numerical values of the congruence field $\mathfrak{E}_{2,p}$, including $\infty$. We thereby gain the advantage that the equation
\weq{2}{\xi = \frac{a\xi+b}{c\xi+d}}
%CONT>>
